% ── Experiment 1: Paragraph insert ─────────────────────────────────────────
% INSERT AT END OF \paragraph{Format Collapse as the Primary Extraction Bottleneck}
% IN \subsection{Results: Keyword Extraction}

To empirically separate the entity-knowledge component of Format Collapse
from its syntactic-compliance component, we re-evaluate Task~2 using
Ollama's native JSON-Schema grammar enforcement (\texttt{format} parameter;
Ollama $\geq$0.1.24), which constrains the decoder's sampling distribution
to outputs that are valid instances of the
$\{\texttt{Symptoms},\texttt{Diagnostics},\texttt{Pathogens}\}$ schema
by construction, bypassing Format Collapse at the decoding layer
(Table~\ref{tab:keywords_constrained}).
The constrained No-RAG baseline achieves $\mu$-F$_1 = 0.087$
($\Delta_\mu = +0.075$ vs.\ native No-RAG).
The large positive $\Delta_{{F_1}}$ on Symptoms (+0.074) confirms that entity knowledge was present in the model's parametric memory but syntactically inaccessible under native generation.
The engineering implication is direct: constrained structured decoding is
a \emph{mandatory architectural prerequisite} for 4-bit local keyword
extraction, not an optimisation; and this experiment establishes the knowledge
ceiling the bare parametric model can achieve once the compliance bottleneck
is removed.
